\section{Exact period ranks for specified Nielsen components}
\label{sec:nielsen-rank-algorithm}

\paragraph{Scope and check status.}
This section gives a terminating exact algorithm, not a practical running-time
bound or a table for unbounded coefficient degree.  It constructs an algebraic
generic cover of the specified component and computes its period rank by
linear algebra in a finitely generated field.  No period integral is used.
The proof and the finite example below are new in this run.  The parent
completed an independent whole-proof review of D-NR-01--04 on
21 November 2026, including the primary Romagny--Wewers inputs, and
recorded PASS.  The finite matrix example was executed exactly; the
general generic-model search has not been executed.  The classical algebraic algorithms used are
exact arithmetic in algebraic numbers and finitely generated fields,
Gr\"obner-basis elimination, and isolation of algebraic roots.

\subsection*{D-NR-01. Input and an exact monodromy subroutine}

The input consists of a finite group $G$, an ordered generating tuple
$\boldsymbol g=(g_1,\ldots,g_r)$ of nonidentity elements with product one,
and an absolutely irreducible representation $V$ over an explicitly given
imaginary quadratic field $K$.  Its dual is not isomorphic to $V$, and its
two multiplicities in $H^0(C,\omega_C)$ are assumed to be $(n,n)$.
The desired output is the generic rank on the connected ordered Hurwitz
component containing $\boldsymbol g$, modulo simultaneous inner conjugacy.
Finite additional level structures do not change this rank.  Cases $r\le3$
have rank zero.  We henceforth take $r\ge4$.  All applications here have
$g(C)\ge2$.

We first describe a subroutine which determines the labelled monodromy of
an algebraic $G$-cover over $\overline{\mathbf Q}$.  Choose a primitive
function $y$ for its function field over $\overline{\mathbf Q}(z)$, with
separable equation
\[
                 p(z,y)=0,\qquad \deg_y p=|G|.                 \tag{NR.1}
\]
Clear denominators so that $p$ is polynomial in both variables; this adds
the zeros of the cleared denominators to the exceptional set.
The deck maps are part of the exact input.  Include among the exceptional
points the poles of coefficients, the zeros of the leading coefficient and
the discriminant of $p$, and the finitely many points needed to evaluate
the chosen rational deck-map formulas.  These are an effectively computable
finite set of algebraic points.  Standard polygon generators of the
punctured sphere can be perturbed, using rational vertices, to avoid this
set while retaining their homotopy classes relative to the actual branch
points.  Small disjoint rational boxes about the exceptional points make
this a finite planar construction.  A loop around an extra discriminant
point has trivial monodromy on the normalization; such a point is not
silently treated as a new branch point.

Here is an entirely algebraic decision procedure for lifting one polygon
edge.  Subdivide it into $2^a$ equal segments.  At every subdivision vertex,
isolate the $|G|$ distinct algebraic roots of (NR.1).  Let $\delta$ be the
minimum distance between these roots at that vertex.  Choose rational
centres within $\delta/128$ of the roots and a common rational radius in
$[\delta/32,\delta/16]$.  Isolation and comparisons of algebraic real
numbers produce these choices exactly.  The resulting closed discs are
pairwise disjoint and each initially contains exactly one root.

Write $p(z,w)=\sum_{j=0}^N a_j(z)w^j$ and
$a_j(z)=\sum_k a_{jk}z^k$, where $N=|G|$.
For the segment $[z_0,z_1]$, put $h=|z_1-z_0|$ and choose a rational
$Z\ge\max(|z_0|,|z_1|,1)$.  For a root disc with centre $c_i$ and
radius $s_i$, choose a rational $R_i\ge|c_i|+s_i$ and a rational
$\varepsilon_i>|c_i-\alpha_i|$, still smaller than $\delta/128$.
On its boundary, factorization at the algebraic endpoint gives the
strictly positive lower bound
\[
 \ell_i=|a_N(z_0)|(s_i-\varepsilon_i)
               \prod_{j\ne i}(|c_i-\alpha_j|-s_i)>0.
\]
Choose a rational $L_i$ with $\ell_i/2<L_i<\ell_i$.
The coefficient derivative estimate on the straight segment gives
\[
 |p(z,w)-p(z_0,w)|\le h D_i,\qquad
 D_i\ge\sum_{j=0}^N R_i^j\sum_{k\ge1}k|a_{jk}|Z^{k-1},
 \qquad h^2D_i^2<L_i^2,                                    \tag{NR.2}
\]
where $D_i$ is a rational upper bound, chosen within a factor two after
adding one if the exact expression is zero.  All comparisons and bounds
in this test are computable with algebraic numbers and rational intervals;
$h^2$ is rational because the polygon vertices are rational.  Test the
last strict inequality for every root disc.  If it passes, the polynomial
homotopy $p(z_0,w)+u(p(z,w)-p(z_0,w))$, $0\le u\le1$, has no zero
on the boundary.  Its winding number, hence its number of zeros in the
disc, is constant.  This is the elementary Rouch\'e argument, and it proves
that exactly one root remains in each disc throughout the segment.
Identify the endpoint root by exact disc membership.  If any inequality
fails, increase $a$ and start the edge again.

This process terminates.  On the compact edge the nonvanishing leading
coefficient has a positive minimum, the roots are bounded, and their
minimum separation is positive.  The prescribed radii and centres make
every factor in $\ell_i$ uniformly bounded below by a positive constant:
$s_i-\varepsilon_i\ge3\delta/128$ and
$|c_i-\alpha_j|-s_i>\delta/2$ for $j\ne i$.
The chosen rational $L_i$ retain at least half this lower bound.
The $R_i$, $Z$ and hence the $D_i$ are uniformly bounded above; take their
rational upper bounds within a fixed additive one when choosing them.
Dyadic subdivision sends $h$ to zero, so eventually every inequality
(NR.2) holds.  All accepted steps are finite exact rational certificates.
This proves both the lifting result and its error control without any
uncertified analytic continuation.

At the common base point choose one point of the fibre.  Its $G$-orbit
identifies the fibre with the regular $G$-set.  Monodromy commutes with the
deck action, so the computed permutations recover elements of $G$, unique
up to simultaneous inner conjugacy.  At increasing real branch values
use the standard ordered positive polygon generators, transported from
$0,1,\ldots,r-1$ through the increasing real chamber.  The same tails,
loop orientations and multiplication order are used for the input tuple.
With a left deck action, continuation acts by right multiplication on
the fibre labelled by $G$; if function pullbacks instead encode the
opposite action, invert that encoding before the tuple comparison.
In this convention the standard positive generators have ordered product
one.  This gives a
terminating exact test for equality with the input inner Nielsen tuple.

\subsection*{D-NR-02. Constructing the specified generic component}

Put $k=\overline{\mathbf Q}$ with its standard computable presentation by
minimal polynomials and isolating rectangles.  Normalize the first three
ordered branch points to $0,1,2$, and let
\[
 B=\{(t_4,\ldots,t_r):0,1,2,t_4,\ldots,t_r
                  \text{ are distinct}\}\subset\mathbf A^{r-3}_k.
                                                               \tag{NR.3}
\]
This is a dense open coordinate chart in the ordered moduli of punctured
spheres.  Write $K_0=k(t_4,\ldots,t_r)$.  The following search is an actual
enumeration of finite algebraic descriptions.

All enumerations below are dovetailed: enumerate complete finite packages
in increasing total description length, including the defining polynomial
of the coefficient field, every chart, map, and certificate.  At each length
there are only finitely many strings in a fixed effective alphabet.  The
phrase ``for each $F$'' therefore never means finishing an infinite search
over the first field before proceeding to the next.  Since every individual
test terminates, every valid finite package is eventually tested.

Enumerate monic irreducible polynomials $P(T)\in K_0[T]$ and put
$F=K_0[T]/(P)$.  Irreducibility here is an effective test: after clearing
denominators and removing the polynomial content in the $t_i$, Gauss's
lemma reduces it to factorization of a multivariate polynomial over the
computable algebraically closed field $k$.  Finite-degree coefficient
equations with nonzero leading coefficients decide that factorization.
For each $F$, enumerate projective curves, rational maps
and their finite affine-chart descriptions over $F$.  Retain a package only
when it passes all the following exact tests: the curve is smooth and
geometrically connected; its map $\pi:C\to\mathbf P^1_F$ is finite of
degree $|G|$; the specified maps give a faithful action of $G$ over
$\mathbf P^1_F$; and the branch locus is exactly the ordered divisor
\[
                    0+1+2+t_4+\cdots+t_r.                    \tag{NR.4}
\]
We may restrict to the genus prescribed by Riemann--Hurwitz for the input
tuple.  Include a primitive plane equation for $F(C)/F(z)$ and explicit
birational identifications in the package, so the monodromy subroutine has
its required presentation.

For completeness, the tests just listed are algebraic decisions, not
oracles for the desired period rank.  Polynomial identities and equality
of rational maps are ideal-membership tests.  Smoothness is the Jacobian
criterion on finitely many affine charts.  Dimensions, degrees and the
Hilbert polynomial are computed from initial monomial ideals; the constant
term of the curve's Hilbert polynomial gives its genus.  Geometric
irreducibility can be tested by absolute decomposition.  Alternatively,
in the enumerated package one can require a birational identification with
an absolutely irreducible plane equation, a smooth one-dimensional
projective curve, and a zero-dimensional complement to that identification.
Absolute irreducibility of a plane equation is decided by trying the
finitely many possible positive factor degrees and testing the coefficient
equations over an algebraic closure.  Enumerate the possible nonzero
leading coefficients of each factor, fix one scaling coefficient to one,
and impose each required nonvanishing by an auxiliary equation $uv=1$.
Gr\"obner elimination decides solvability of each such finite system.
This alternative avoids a general primeness
oracle.  The absence of base points and the identities on overlaps verify
the chart descriptions of morphisms.  A nonconstant morphism from the
smooth projective curve to $\mathbf P^1$ is finite; its degree is checked
in its primitive function-field presentation.  The group law and
distinctness of all $|G|$ automorphisms then show that it is a $G$-cover.
Finally the torsion module $\Omega^1_{C/\mathbf P^1}$ is computed from
the differentials of the chart equations.  Its support, projected to
$\mathbf P^1$, is computed by elimination and is compared with (NR.4),
including the chart at infinity.  These operations terminate over the
computable characteristic-zero field $F$.

Spread a retained package to a good family over a nonempty open in the
normalization of $B$ in $F$.  This step is effective as well: clear all
denominators and remove the images of the nonsmooth, nonfinite,
nonflat, undefined-map and wrong-branch loci.  Jacobian ideals,
Fitting ideals, affine-cover certificates and elimination compute these
loci.  Generic flatness can be effected by inverting the finitely many
leading coefficients in the relevant module Gr\"obner bases.  Generic
geometric connectedness and proper smoothness then give geometric
connectedness throughout this open.  Each excluded locus is proper,
because the generic package has passed the tests.  Equivalently one may
enumerate base localizations and their finite Jacobian, Fitting and
covering certificates; generic validity guarantees that a successful
localization occurs, and the tests themselves terminate.

We need an open of the \emph{base} over which \emph{every} sheet is good.
After clearing coefficients of $P$ and inverting its discriminant, the
monogenic algebra is already finite etale over an open of $B$.  If a
further nonzero element of this algebra has to be inverted, invert its
nonzero field norm down to $K_0$.  After finitely many such steps we obtain
explicit $D\ne0$ and
\[
 U=B\cap\{D\ne0\},\qquad
 S=\operatorname{Spec}\bigl(k[t_4,\ldots,t_r,D^{-1},T]/(P)\bigr)
       \longrightarrow U                                    \tag{NR.5}
\]
finite etale, with the entire package a smooth family of connected
$G$-covers over $S$.  Factors needed to clear coefficients and branch
diagonals are included in $D$.  Because $P$ is irreducible over $K_0$,
$S$ is integral and $S(\mathbf C)$ is connected.

Choose integers $2<t_4<\cdots<t_r$ with $D\ne0$ by enumeration.
Such integers exist: that ordered integer set is Zariski dense, as follows
by induction on the number of variables from the finiteness of the zeros
of a nonzero univariate polynomial.  Compute \emph{every} root of
$P$ at these integers.  Each is an algebraic point of $S$ and gives an
exact algebraic $G$-cover with real ordered branch points.  Apply D-NR-01
to each of these finitely many covers.  If one has the input tuple up to
simultaneous inner conjugacy, retain this generic package.  Otherwise
continue the enumeration.

We prove both the correctness and termination of this selection.
The inner Hurwitz space of complex connected $G$-covers is a finite
topological covering of branch configuration space; its fibre is exactly
the finite set of generating product-one tuples modulo inner conjugacy.
The local charts and the analytic classifying map for a family are proved
in \cite[Theorem~3.1, Propositions~3.2 and~3.4, pp.~319--321]{NR:RW}.
Ordered branch points and the normalization (NR.3) are obtained by base
change and do not alter this assertion.  The classifying map from
$S(\mathbf C)$ is a map of finite covering spaces over $U(\mathbf C)$.
It is therefore a covering onto one connected Hurwitz component over
$U$: locally it maps a sheet isomorphically to a sheet, so its image is
open; finiteness over $U$ makes the image closed.  Consequently its finite
fibre over the chosen integer branch configuration contains every inner
Nielsen point in that component.  This is why all roots of $P$, not one
arbitrary root, are tested.

Restriction from $B$ to the nonempty Zariski open $U$ does not split a
connected smooth Hurwitz component: an irreducible complex algebraic
variety remains irreducible on a nonempty Zariski open.  The algebraic
Hurwitz scheme and its agreement with this analytic covering are
\cite[Theorem~4.11, pp.~328--331, and Remark~4.15(i)]{NR:RW}.
Thus a successful tuple test identifies exactly the required component.
Since $S\to U$ is finite etale, its family dominates that component and
has the same generic period rank.

\paragraph{Grothendieck's algebraic framework.}
This construction sits in the algebraic theory of tame Hurwitz schemes
over $\mathbf Z[1/|G|]$: after passage to characteristic zero, the
branch-coordinate morphism is a finite etale cover of the appropriate
configuration space.  The construction and comparison used here are
\cite[Theorem~4.11 and Remark~4.15]{NR:RW}, together with their
finite-field-of-definition Theorem~4.3.  They permit the passage from a
topologically specified cover to a finite algebraic generic description.
The algorithm itself uses only characteristic zero and computes an exact
generic Hodge matrix over a finite extension of
$\overline{\mathbf Q}(t_4,\ldots,t_r)$.  No stronger effective theorem in
positive characteristic is being asserted.  In particular we do not
attribute a direct field-of-definition statement to an unspecified SGA
locator; Romagny--Wewers explicitly discuss that source issue.

Conversely let $\mathcal H$ be the component specified by the input.
It is defined over $k$, and its function field is a finite extension of
$K_0$.  A cover representing its geometric generic point is defined over
a finite further extension $F$ of that function field: its equations,
the cover map and the finitely many labelled group automorphisms use only
finitely many algebraic coefficients.  This also follows directly from
the finite-field-of-definition theorem
\cite[Theorem~4.3 and its proof, pp.~323--327]{NR:RW}.
The center of $G$ is allowed; we do not assume that the coarse generic
point itself carries a universal cover.  The finite further extension
removes exactly this possible field-of-moduli obstruction.
Characteristic zero supplies a primitive element for $F/K_0$.
The resulting cover, any smooth projective embedding and its finite map
descriptions occur in the enumeration.  They spread to a package of the
form (NR.5), and its tuple test succeeds.  Every earlier candidate is
decided in finitely many steps.  Therefore the search terminates.

\subsection*{D-NR-03. A finite matrix whose rank is the answer}

For the accepted generic cover set
\[
 A=\operatorname{Hom}_{G}(V_F,H^0(C,\omega_C)),\qquad
 A'=\operatorname{Hom}_{G}(V_F^*,H^0(C,\omega_C)).
                                                               \tag{NR.6}
\]
Both spaces have dimension $n$.  Choose a basis $e_1,\ldots,e_m$ of $V$
and its dual basis.  Define
\[
 \mu:A\otimes_F A'\longrightarrow H^0(C,\omega_C^2)^G,
 \qquad
 \mu(a\otimes b)=\sum_{j=1}^m a(e_j)b(e_j^*).                \tag{NR.7}
\]
The expression is independent of the dual bases and is $G$-invariant.

All parts of (NR.6)--(NR.7) can be computed exactly.  One explicit
terminating method to compute $H^0(C,\omega_C)$ is useful here.  Every
rational differential has the form $f\,dz$ with $f\in F(C)$.
Enumerate these differentials.  On a finite smooth affine cover, regularity
is membership in the finitely presented K\"ahler differential module.
After a common denominator is taken, membership of $v/b$ is the module
membership test $v\in b\Omega^1$; the differential module is torsion-free
because the curve is smooth.  Module Gr\"obner bases decide it.  Keep
$F$-linearly independent regular differentials until $g(C)$ are found.
Riemann--Roch guarantees termination.  Linear dependence over $F$ is
decided by writing elements in the basis of $F(C)/F(z)$, clearing their
$z$-denominators, and comparing polynomial coefficients.  Pullback by each
deck map is exact in the same function field.  The intertwiner equations
defining (NR.6) are consequently finite linear equations over $F$.

Write $b_1=0,b_2=1,b_3=2,b_i=t_i$ and
$Q(z)=\prod_{i=1}^r(z-b_i)$.  Invariant quadratic differentials descend as
\[
 H^0(C,\omega_C^2)^G
  =\pi^*H^0\bigl(\mathbf P^1_F,\omega_{\mathbf P^1}^2(B)\bigr),
 \qquad
 \frac{z^j(dz)^2}{Q(z)},\quad0\le j\le r-4,                 \tag{NR.8}
\]
where equality means the canonical pullback identification.  At a branch
point $z=u^e$, invariance of a holomorphic quadratic differential makes
its coefficient a multiple of $u^{e-2}$ modulo powers of $u^e$; descent
therefore allows exactly a simple pole.  At an unramified point it is
holomorphic.  At infinity the degree bound in (NR.8) is precisely the
condition of regularity.

If $\mu(a\otimes b)=f(dz)^2$, its descended coefficient can also be
computed as $|G|^{-1}\operatorname{Tr}_{F(C)/F(z)}(f)$.
The trace is an exact finite matrix trace.  Express these outputs in the
basis (NR.8) to obtain an explicit
\[
                       (r-3)\times n^2\text{ matrix }M
                       \text{ over }F.                      \tag{NR.9}
\]

We prove that its rank is the desired geometric rank.  For a smooth curve,
a Cech representative $\{\theta_{ij}\}$ of a first-order deformation
sends a holomorphic differential $\alpha$ to the class
$\{\theta_{ij}\mathbin{\lrcorner}\alpha\}\in H^1(C,\mathcal O_C)$.
Serre duality pairs this class with $\beta$ by pairing the deformation
with $\alpha\beta$.  This derives the multiplication codifferential of
the Jacobian period map.  For a $G$-invariant deformation only the group
average of a product contributes.  On $V\otimes V^*$, the group average
of a rank-one matrix unit is $m^{-1}$ times the identity, by Schur's
lemma.  Restriction to the primitive $K$-factor thus gives (NR.7), up to
a fixed nonzero scalar coming from the normalization of its polarization.
The cotangent space of the signature-$(n,n)$ unitary period domain is
$A\otimes A'$.  Polarization identifies the conjugate filtration with
the annihilator of the first; no extra independent period directions are
omitted in this description.

The tangent space of the Hurwitz component is
$H^1(C,T_C)^G$.  The local computation $z=u^e$ gives
$(\pi_*T_C)^G=T_{\mathbf P^1}(-B)$: an invariant vector field is
$u h(u^e)\partial_u$ and descends to $e z h(z)\partial_z$.
Taking cohomology and using exactness of finite-group invariants gives
\[
 H^1(C,T_C)^G=H^1(\mathbf P^1,T_{\mathbf P^1}(-B)),
 \qquad \dim=r-3.                                           \tag{NR.10}
\]
The branch-coordinate map is etale, so the family retained in D-NR-02
has this entire tangent space.  Duality identifies the dual of (NR.10)
with (NR.8).  Consequently (NR.9) is the full period codifferential on
the specified component, not a restriction to a conveniently chosen
subfamily.

Gaussian elimination in $F$ now outputs $q=\operatorname{rank}_F M$.
Field equality is decidable by reduction modulo $P$ and equality of
rational functions in the algebraically independent $t_i$.
One nonzero $q$-minor and a row-reduction identity with exactly $q$
nonzero rows are finite exact lower- and upper-rank certificates.
Base change to $\mathbf C$ preserves rank.  On a nonempty open of the
Hurwitz component this equals the complex period differential rank,
and hence the dimension of its image.  This proves a terminating exact
algorithm for every specified input tuple satisfying the hypotheses.

\subsection*{D-NR-04. The smaller trace formula for natural imprimitive groups}

For the natural representation of $G(d,e,m)$, use the actual intermediate
tower already constructed in D-IM-01:
\[
             Y=C/\ker\chi\ \xrightarrow{\ q\ }\quad
             X=C/H\ \xrightarrow{\ f\ }\ \mathbf P^1,
 \qquad V=\operatorname{Ind}_H^G\chi.                         \tag{NR.11}
\]
Choose bases $\alpha_1,\ldots,\alpha_n$ and
$\beta_1,\ldots,\beta_n$ of the two reciprocal eigenspaces of
$H^0(Y,\omega_Y)$.  Their products are invariant under the cyclic group,
so each descends to a meromorphic quadratic differential
$\gamma_{ab}$ on $X$, with at most simple poles at the cyclic branch
points.  Equivalently,
\[
 \gamma_{ab}=\frac1d\operatorname{Tr}^{(2)}_{Y/X}
                          (\alpha_a\beta_b).                \tag{NR.12}
\]
Here the superscript means trace of quadratic differentials; in a common
base coordinate it is the ordinary field trace of their coefficients.
Frobenius reciprocity identifies the coset basis of the induced module
with the conjugate pullbacks of these eigendifferentials.  The dual-basis
contraction in (NR.7) sums exactly over the $m$ cosets.  Therefore it is,
up to a fixed nonzero scalar,
\[
                     \operatorname{Tr}^{(2)}_{X/\mathbf P^1}
                                      (\gamma_{ab}).         \tag{NR.13}
\]
This also proves the formula directly at the level of variations and
their polarization pairings; it is not an equality inferred from their
dimensions.  Local trace at a ramification point, or (NR.8), shows that
(NR.13) has at most simple poles at the original branch divisor $B$.

Thus the exact rank of any specified natural-representation Nielsen
component is
\[
 \boxed{\dim_F\operatorname{span}\left\{
       \operatorname{Tr}^{(2)}_{X/\mathbf P^1}(\gamma_{ab}):
                       1\le a,b\le n\right\}.}              \tag{NR.14}
\]
The finite algebraic model search above makes this a computable formula,
including for arbitrary ramification of $X\to\mathbf P^1$.
It retains that map and imposes no submersion assumption when forgetting
it.  The universal bounds and the ranks $15$, $14$ and $1$ proved in
D-IM-03--08 remain distinct claims, not a universal replacement of
(NR.14) by one of the maxima.

For each fixed $m,d,e,n$, this also gives a finite procedure to tabulate
all answers.  Faithfulness and Chevalley--Weil give
$\sum_j\operatorname{rank}(1-g_j)=2m+2n$, so $r\le2m+2n$.
There are finitely many tuples in the finite group for these lengths.
Product one, generation, and signature are finite exact tests; run the
algorithm on each passing tuple.  Computing finite braid orbits can remove
duplicate work, but is not needed for termination.  This is not a claim
that the resulting table for every unbounded $m$ has been executed.

\paragraph{A finite check of the matrix step.}
For a cyclic cubic cover of $\mathbf P^1$ with six branch exponents one
and six branch exponents two, write
\[
 y^3=P(z)Q(z)^2,\qquad \deg P=\deg Q=6,
\]
with twelve distinct roots.  The two eigenspace bases are
$z^a dz/y$ and $z^bQ(z)dz/y^2$, $0\le a,b\le4$.
Their products are $z^{a+b}(dz)^2/(PQ)$, and the exponents $a+b$
span $0,\ldots,8$.  The exact matrix has rank nine.  The accompanying
small script checks all local valuations, the two dimensions five, and
an identity $9$-minor.  This illustrates and certifies the matrix step;
it is not an execution of the generic-model search for a large
$G(d,e,m)$, and it makes no additional claim about those components.
